\documentclass[10pt,a4paper]{amsart}
\usepackage[margin=19mm]{geometry}
\usepackage{amsmath,amssymb,mathtools}
\usepackage[scr=rsfs,cal=euler]{mathalfa}
\usepackage{xfrac}
\usepackage[unicode,hidelinks,bookmarks=false]{hyperref}
\newcommand{\ie}{i.e.\ }
\newcommand{\cf}{cf.\ }
\newcommand{\resp}{respectively\ }
\newcommand{\Subsection}{\subsection}
\providecommand{\nobreakdash}{\nobreak\hbox{-}}
\setlength{\emergencystretch}{2em}
\multlinegap0pt
\usepackage{amssymb}
\usepackage{tikz}
\newtheorem{thm}{Theorem}
\newtheorem{lemma}{Lemma}
\newcommand{\CC}{\mathbb{C}}
\newcommand{\PP}{\mathbb{P}}
\title{Type \(C\) Richardson Boundaries and Newton--Okounkov Degenerations of Odd-Dimensional Projective Space}
\author{Zhaoyang Liu}
\date{Source: arXiv:2608.11109v1 (CC BY 4.0)}
\begin{document}
\maketitle

\noindent\emph{Source and licence.} Type \(C\) Richardson Boundaries and Newton--Okounkov Degenerations of Odd-Dimensional Projective Space. Version arXiv:2608.11109v1 (CC BY 4.0), distributed by arXiv under the Creative Commons Attribution 4.0 International licence, http://creativecommons.org/licenses/by/4.0/, as recorded in the arXiv licence metadata for this exact version and shown on the abstract page https://arxiv.org/abs/2608.11109v1. Full licence notice: \texttt{licenses/arxiv-2608.11109-CC-BY-4.0.txt}.\par\smallskip

\noindent\textit{Editorial scope.} Counted unit: the article's thm thm:richardson-boundary (source0.tex lines 450--469). The article's complete original proof of it is delivered with it (source0.tex lines 553--761, one \texttt{proof} environment of 209 lines, which the article places away from the statement). No mathematics is written, repaired or re-derived: the statement and the proof are the article's own lines, in the article's own notation, and no retained line is substituted or re-flowed.  Retained uncounted as the source-local context the proof needs: lines 472--483; lines 503--512; lines 514--524.  Nothing was omitted from any retained span, so the package carries no omission record.  No retained line contains a citation, so the wrapper carries no bibliography and no bibliography entry.  No retained line contains a figure environment, an image-inclusion command or drawing code, so no image is delivered and the package declares no asset dependency.  Editorial formatting only, in addition: an AMS \texttt{amsart} preamble that restates the article's own declarations and macros used by the retained lines, namely the environments \texttt{thm}, \texttt{lemma} on separate counters, together with the standard packages those lines need.  The retained environments print in this document as Theorem 1, Lemma 1 in order of appearance, whereas the article prints the same items with the numbering of its own full document.  The complete native source of the article as distributed in the arXiv e-print for this exact version is bundled unchanged as \texttt{sources/207388/source0.tex}.
\begin{lemma}\label{lem:southeast-elimination}
Let \(A\in GL(V)\) where \(V\) is an \((N+1)\)-dimensional vector space with ordered basis \(e_0,\dots,e_N\).
Define the southeast principal minors
\[
\Delta_k(A)=
\det A_{\{k,k+1,\dots,N\},\{k,k+1,\dots,N\}}
\qquad(1\le k\le N),
\]
and set \(\Delta_{N+1}(A)=1\).
Then \(A\) admits a unique Gauss factorization \(A=u^{-1}b\), with \(u\in U_+\) and \(b\in B_-\), if and only if \(\Delta_k(A)\neq 0\) for \(1\le k\le N\).
Moreover, if \(A\in Sp(V,J)\), then both factors \(u\) and \(b\) belong to \(Sp(V,J)\).
\end{lemma}

\begin{lemma}\label{lem:jacobi-complementary-minor}
Let \(A\in GL_{N+1}(\CC)\), with rows and columns indexed by \(0,\ldots,N\). Let \(I,J\subset\{0,\ldots,N\}\) have the same cardinality, and write \(I^c,J^c\) for their complements, each in increasing order. Then
\[
\det A_{I,J}
=
(-1)^{\sum_{i\in I}i+\sum_{j\in J}j}
\det(A)\,
\det (A^{-1})_{J^c,I^c}.
\]
\end{lemma}

\begin{lemma}\label{lem:symplectic-inverse-last-column}
Let \(v\in U_+\cap Sp(V,J)\), and write its first row as
\[
q=e_0^tv=(q_0,q_1,\ldots,q_N).
\]
Then the last column of \(v^{-1}\) is obtained from the first row of \(v\) by reversing the order and alternating signs:
\[
(v^{-1})_{iN}=(-1)^{i+1}q_{N-i}=(-1)^{i-N}q_{N-i}
\qquad(0\le i\le N).
\]
\end{lemma}

\begin{thm}\label{thm:richardson-boundary}
The type \(C_n\) open locus in \(\PP^N=\PP^{2n-1}\) is
\[
\eta(S)=
\{[p_0:\cdots:p_N]\in \PP^N:
 p_0p_NF_1\cdots F_{n-1}\ne0\}.
\]
Consequently the reduced type \(C_n\) boundary is
\[
D_C
=
\operatorname{div}\bigl(p_0p_NF_1\cdots F_{n-1}\bigr)_{\mathrm{red}}
=
\{p_0=0\}
\cup
\left(\bigcup_{r=1}^{n-1}\{F_r=0\}\right)
\cup
\{p_N=0\}.
\]
\end{thm}

\begin{proof}[Proof of Theorem~\ref{thm:richardson-boundary}]
By definition, \(\eta:S\longrightarrow \PP^N,b\longmapsto [e_N^tb\dot w_0]\). Since \(S=B_-\cap U_+MU_+\), every \(b\in S\) can be written as \(b=uMv\) for some \(u,v\in U_+\cap Sp(V,J)\). Let \(q=(q_0,\ldots,q_N):=e_0^tv\) be the first row of \(v\). Since \(v\in U_+\), one has \(q_0=1\). The last row of \(u\) is \(e_N^t\), and the last row of \(M\) is \(e_0^t\). Therefore
\[
e_N^tb=e_N^tuMv=e_N^tMv=e_0^tv=q.
\]
The homogeneous coordinates of \(\eta(b)\) are \(p=qJ\), namely
\[
p_i=(qJ)_i=(-1)^{N-i}q_{N-i}.
\]
Thus \(p_i=(-1)^{N-i}q_{N-i}\). Since \(N=2n-1\) is odd,
\[
p_N=q_0=1,
\qquad
p_0=-q_N.
\]
Moreover, substituting \(p_i=(-1)^{N-i}q_{N-i}\) gives
\[
F_r(p)=-F_r(q)
\qquad
(1\le r\le n-1).
\]
Thus the stated open set in \(p\)-coordinates is the image, under \(q\mapsto [qJ]\), of the \(q\)-coordinate condition
\[
q_0=1,\qquad F_0(q)F_1(q)\cdots F_{n-1}(q)\ne0.
\]
It remains to compare this condition with the Gauss factorization condition for \(Mv\).

For an invertible matrix \(A\), write
\[
\Delta_k(A)=
\det A_{\{k,k+1,\ldots,N\},\{k,k+1,\ldots,N\}},
\qquad
1\le k\le N.
\]
By Lemma~\ref{lem:southeast-elimination}, \(A\) admits a unique Gauss factorization \(A=u^{-1}b\) with \(u\in U_+\) and \(b\in B_-\) if and only if
\[
\Delta_k(A)\ne0\qquad(1\le k\le N).
\]
Now take \(A=Mv\). The rows of \(A\) are
\[ A_0=-v_N,
\quad
A_i=-v_i\quad(1\le i\le N-1),
\quad
A_N=v_0,
\]
where \(v_i\) is the \(i\)-th row of \(v\). Since \(v\in U_+\), all diagonal entries of \(v\) are \(1\) and all entries below the diagonal are \(0\). Thus \(A=Mv\) has the corresponding form. For example, when \(N=5\),
\[
A=
\begin{pmatrix}
0&0&0&0&0&-1\\
0&-1&*&*&*&*\\
0&0&-1&*&*&*\\
0&0&0&-1&*&*\\
0&0&0&0&-1&*\\
1&q_1&q_2&q_3&q_4&q_5
\end{pmatrix}.
\]
For \(1\le k\le N-1\), expand \(\Delta_k(A)\) along its first column, i.e. the global column \(k\). In that column the only possible nonzero entries are \(-1\) in row \(k\) and \(q_k\) in row \(N\). Hence
\[
\Delta_k(A)
=
-\Delta_{k+1}(A)
+(-1)^{N-k}q_kC_k,
\]
where
\[
C_k=
\det A_{\{k,\ldots,N-1\},\{k+1,\ldots,N\}}.
\]
To compute \(C_k\), apply Lemma~\ref{lem:jacobi-complementary-minor} to the matrix \(A\). Put
\[
I=\{k,\ldots,N-1\},
\qquad
J=\{k+1,\ldots,N\}.
\]
Then
\[
I^c=\{0,\ldots,k-1,N\},
\qquad
J^c=\{0,\ldots,k\}.
\]
Since \(A\in Sp(V,J)\), \(\det(A)=1\), and the sign in Jacobi's identity is \((-1)^{N-k}\). Thus
\[
C_k
=
(-1)^{N-k}
\det(A^{-1})_{J^c,I^c}.
\]
Here the complementary row set in \(A\) is \(I^c=\{0,\ldots,k-1,N\}\), and the complementary column set is \(J^c=\{0,\ldots,k\}\); in Jacobi's formula they occur as rows \(J^c\) and columns \(I^c\) of \(A^{-1}\). Since \(A=Mv\), we have \(A^{-1}=v^{-1}M^{-1}\). The inverse of \(M\) is again a signed permutation matrix:
\[
(M^{-1})_{0N}=1,
\qquad
(M^{-1})_{N0}=-1,
\qquad
(M^{-1})_{ii}=-1\quad(1\le i\le N-1),
\]
with all other entries zero. Thus, when \(N=5\), the shapes of \(v^{-1}\) and \(M^{-1}\) are
\[
\begin{array}{c@{\qquad}c}
v^{-1}=
\begin{pmatrix}
1&*&*&*&*&-q_5\\
0&1&*&*&*&q_4\\
0&0&1&*&*&-q_3\\
0&0&0&1&*&q_2\\
0&0&0&0&1&-q_1\\
0&0&0&0&0&1
\end{pmatrix}
&
M^{-1}=
\begin{pmatrix}
0&0&0&0&0&1\\
0&-1&0&0&0&0\\
0&0&-1&0&0&0\\
0&0&0&-1&0&0\\
0&0&0&0&-1&0\\
-1&0&0&0&0&0
\end{pmatrix}.
\end{array}
\]
The displayed last column of \(v^{-1}\) is given by Lemma~\ref{lem:symplectic-inverse-last-column}.
Right multiplication by \(M^{-1}\) sends column \(0\) to minus the \(N\)-th column, sends column \(i\) to minus the \(i\)-th column for \(1\le i\le N-1\), and sends column \(N\) to the \(0\)-th column. Hence \(A^{-1}=v^{-1}M^{-1}\) has the following form when \(N=5\):
\[
A^{-1}=
\begin{pmatrix}
q_5&*&*&*&*&1\\
-q_4&-1&*&*&*&0\\
q_3&0&-1&*&*&0\\
-q_2&0&0&-1&*&0\\
q_1&0&0&0&-1&0\\
-1&0&0&0&0&0
\end{pmatrix}.
\]
For the required minor, this shape shows that the matrix is triangular except for the column coming from the \(N\)-th column of \(v^{-1}\):
\[
\det(A^{-1})_{J^c,I^c}=q_{N-k}.
\]
Indeed, the column set \(I^c\) consists of the columns \(0,\ldots,k-1,N\); after multiplying by \(M^{-1}\), these are \(-\)the \(N,1,\ldots,k-1\) columns of \(v^{-1}\), followed by the \(0\)-th column of \(v^{-1}\). The \(0\)-th column is \(e_0\), the columns \(1,\ldots,k-1\) form a unit triangular block, and the remaining entry is
\[
(v^{-1})_{kN}=(-1)^{k-N}q_{N-k};
\]
the triangular expansion gives
\[
\det(A^{-1})_{J^c,I^c}
=(-1)^{k+1}(v^{-1})_{kN}
=q_{N-k},
\]
because \(N\) is odd. Hence
\[
C_k=(-1)^{N-k}q_{N-k},
\]
and therefore
\[
\Delta_k(A)+\Delta_{k+1}(A)=q_kq_{N-k}
\qquad(1\le k\le N-1).
\]

Now \(\Delta_N(A)=A_{NN}=q_N=F_0(q)\). Starting from this value and applying the recurrence backwards gives
\[
\Delta_{N-s}(A)=H_s(q)
\qquad(0\le s\le N-1),
\]
where
\[
H_s(q)=\sum_{j=0}^s(-1)^jq_{s-j}q_{N-s+j}.
\]
For \(0\le s\le n-1\), this is \(F_s(q)\). For \(s\ge n\), the middle terms in the alternating sum cancel in pairs because \(N=2n-1\) is odd, and \(H_s(q)=F_{N-1-s}(q)\). Thus every southeast principal minor \(\Delta_k(A)=\Delta_k(Mv)\), \(1\le k\le N\), is one of \(F_0(q),\ldots,F_{n-1}(q)\), and each of these \(F_r(q)\)'s occurs. Therefore all southeast principal minors of \(Mv\) are nonzero if and only if \(F_0(q),F_1(q),\ldots,F_{n-1}(q)\) are nonzero.

Conversely, take any row vector \(q=(1,q_1,\ldots,q_N)\) satisfying
\[
F_0(q)F_1(q)\cdots F_{n-1}(q)\ne0.
\]
It remains to realize \(q\) as the first row of an element of \(U_+\cap Sp(V,J)\). Write \(a=(q_1,\ldots,q_{N-1})\) and \(z=q_N\). With respect to the decomposition
\[
V=\CC e_0\oplus \operatorname{Span}(e_1,\ldots,e_{N-1})\oplus \CC e_N,
\]
write
\[
J=\begin{pmatrix}
0&0&1\\
0&J'&0\\
-1&0&0
\end{pmatrix}.
\]
Set
\[
v(q)=
\begin{pmatrix}
1&a&z\\
0&I_{N-1}&J'a^t\\
0&0&1
\end{pmatrix}.
\]
Using \(J'^t=-J'\) and \(J'J'=-I\) for the middle symplectic block, block multiplication gives
\[
v(q)^tJv(q)=
\begin{pmatrix}
0&0&1\\
0&J'&0\\
-1&0&0
\end{pmatrix}
=J.
\]
Thus \(v(q)\in U_+\cap Sp(V,J)\), and its first row is \(q\). Since \(F_0(q),F_1(q),\ldots,F_{n-1}(q)\) are nonzero, all southeast principal minors of \(Mv(q)\) are nonzero. Gaussian elimination gives \(u\in U_+\cap Sp(V,J)\) with
\[
uMv(q)\in B_-.
\]
Then \(b=uMv(q)\in S\), and \(e_N^tb\dot w_0=qJ\). Hence every row vector \(q=(1,q_1,\ldots,q_N)\) satisfying \(F_0(q)F_1(q)\cdots F_{n-1}(q)\ne0\) occurs. By the coordinate identification \(p=qJ\) at the beginning of the proof, this is exactly the reverse inclusion in the stated \(p\)-coordinates.
\end{proof}

\end{document}
